%% toc: P1 7 7.4,7.5,7.6,7.7
\section{Invariant arguments and changes of governing logic}
\label{sec:p01-w02-p01-023-invariants-logic}

The proof methods of \cref{sec:writing-foundations-04} name the invariant
maintained at each step as one item of a proof plan, but do not say how an
invariant argument is carried out. The error taxonomy of
\cref{sec:writing-foundations-07} asks whether a rule belongs to the active
calculus, but has no entry for an argument that silently passes from one
logic to a stronger one. This section supplies both. Its first part treats
invariants as a proof method for constructions, iterations and derivations.
Its second part gives a diagnostic for a change of governing logic inside an
argument, with its repairs, for the case in which the stated logic is
intuitionistic propositional logic and the principle used is classical.

All results below are editorial completions with full proofs in the annex.
They use only the definitions of this section, the classical two-valued
semantics of \cref{sec:prop-syntax-semantics-02}, the Kripke semantics of
\cref{sec:intuitionistic-logic-01}, and the soundness and completeness
theorems of \cref{sec:p08-ipl-kripke-completeness}. The metalogic throughout
is classical, as in those sections.

\subsection{Invariants}
\label{sec:p01-w02-p01-023-invariants}

\paragraph{Purpose, inputs and conclusions.}
An invariant argument proves that every state produced by a stepwise process
has a property, without following the process to any particular state. Its
input is a description of the process (its possible starting states and the
steps it may take) and a candidate property. It has three checks:
the property holds at the start; each single step from a state with the
property leads to a state with the property; and the property, together with
whatever else is known when the argument closes, implies the target
statement. Its conclusion is the target statement for every reachable state,
for every terminal state, or for every derivable item, according to the
closing step used. The same pattern governs a construction by stages, an
iteration that repeats an instruction until a condition fails, and a
calculus whose derivable items are generated by rules.

The method has its classical formulation in program verification. Floyd
attaches a proposition to each connection in the flow of control of a
program, imposes on each command a condition guaranteeing that a true
proposition at entry yields a true proposition at exit, and concludes by
induction on the number of commands executed; termination is shown
separately by an entity that each step decreases and that cannot decrease
indefinitely \cite[p.~19]{Floyd1967}. Hoare's rule of iteration D3 states the
loop case: an assertion preserved by the loop body under the loop condition
holds after the loop, together with the negated condition
\cite[p.~578, \S\,3.4]{Hoare1969}. Hoare also records that these rules give
no basis for proving that a program terminates
\cite[p.~579, \S\,4]{Hoare1969}. The handbook's Hoare calculus keeps the two
cases apart, with a partial-correctness while rule and a separate
total-correctness while rule that carries an integer-valued variant
(\cref{sec:p09-program-logics}, \cref{def:p09-hoare-triple}).
The definitions below separate the same three ingredients (invariant, exit
condition, termination measure) without fixing a programming language.

\begin{hbdefinition}{def:p01-w02-p01-023-step-system}{Step system, run, invariant and inductive property}
A \emph{step system} is a triple \(\langle S,I,\leadsto\rangle\) in which
\(S\) is a set of \emph{states}, \(I\subseteq S\) is the set of
\emph{initial states}, and \({\leadsto}\subseteq S\times S\) is the
\emph{step relation}; \(s\leadsto t\) says that one step can lead from
\(s\) to \(t\).
\begin{enumerate}
  \item A \emph{run} of length \(n\in\mathbb N\) is a sequence
        \(s_0,\ldots,s_n\) of states with \(s_0\in I\) and
        \(s_k\leadsto s_{k+1}\) for every \(k<n\). It \emph{starts at}
        \(s_0\) and \emph{ends at} \(s_n\).
  \item A state is \emph{reachable} if some run ends at it.
        \(\operatorname{Reach}\) is the set of reachable states.
  \item A state \(s\) is \emph{terminal} if there is no \(t\) with
        \(s\leadsto t\).
  \item A \emph{property} is a subset \(P\subseteq S\). It is an
        \emph{invariant} if \(\operatorname{Reach}\subseteq P\).
  \item A property \(P\) is \emph{inductive} if
        \begin{enumerate}
          \item[(I1)] \(I\subseteq P\) (\emph{establishment}), and
          \item[(I2)] for all \(s,t\in S\), if \(s\in P\) and
                \(s\leadsto t\), then \(t\in P\) (\emph{preservation}).
        \end{enumerate}
  \item A \emph{variant} for a property \(P\) is a function
        \(\mu:S\to\mathbb N\) such that \(\mu(t)<\mu(s)\) whenever
        \(s\in P\) and \(s\leadsto t\).
\end{enumerate}
A construction by stages is the case in which the states are the possible
stages and \(s\leadsto t\) says that \(t\) is a permitted next stage after
\(s\); an iteration is the case in which the states are the values of the
program variables and \(s\leadsto t\) is one execution of the loop body
from a state satisfying the loop condition.
\end{hbdefinition}

Being an invariant is the property one wants; being inductive is the
property one can check, because (I1) and (I2) each concern a single state or
a single step. The next theorem says that the check suffices and that it is
always available in principle.

\begin{hbtheorem}{thm:p01-w02-p01-023-invariant-principle}{proved=annex:p01-w02-p01-023-invariant-principle}{Invariant principle}
Let \(\langle S,I,\leadsto\rangle\) be a step system.
\begin{enumerate}
  \item Every inductive property is an invariant.
  \item \(\operatorname{Reach}\) is inductive, and
        \(\operatorname{Reach}\subseteq P\) for every inductive property
        \(P\).
  \item A property \(Q\) is an invariant if and only if some inductive
        property \(P\) satisfies \(P\subseteq Q\).
\end{enumerate}
\end{hbtheorem}

Clause~3 is the source of the most common repair: a property that holds of
every reachable state but fails (I2) must be replaced by a stronger
inductive property, and one always exists. A convenient one is usually not
\(\operatorname{Reach}\) itself, whose description is the problem to be
solved, but a property that records exactly the facts about the current
state that the preservation check needs
(\cref{thm:p01-w02-p01-023-non-inductive}).

\begin{hblemma}{thm:p01-w02-p01-023-variant}{proved=annex:p01-w02-p01-023-variant}{Termination by a variant}
Let \(P\) be an inductive property of a step system
\(\langle S,I,\leadsto\rangle\) and let \(\mu\) be a variant for \(P\).
For every \(s_0\in I\):
\begin{enumerate}
  \item every run starting at \(s_0\) has length at most \(\mu(s_0)\); and
  \item some run starting at \(s_0\) ends at a terminal state.
\end{enumerate}
\end{hblemma}

\begin{hbcorollary}{thm:p01-w02-p01-023-exit}{proved=annex:p01-w02-p01-023-exit}{Closing an invariant argument at the exit}
Let \(P\) be an inductive property of a step system
\(\langle S,I,\leadsto\rangle\), let \(T\) be its set of terminal states,
and let \(Q\subseteq S\) satisfy \(P\cap T\subseteq Q\).
\begin{enumerate}
  \item Every reachable terminal state lies in \(Q\) (partial correctness).
  \item If moreover \(P\) has a variant, then every \(s_0\in I\) starts a
        run that ends at a terminal state, and every such state lies in
        \(Q\) (total correctness).
\end{enumerate}
\end{hbcorollary}

Derivations are handled by the same principle with trees in place of
sequences. A \emph{rule system} on a set \(U\) of items is a set
\(\mathcal R\) of \emph{rule instances}, each a pair
\((\langle u_1,\ldots,u_k\rangle,u)\) with \(k\in\mathbb N\) and
\(u_1,\ldots,u_k,u\in U\); side conditions are part of what makes a pair an
instance, and an instance with \(k=0\) is an axiom. A \emph{derivation} is a
finite tree whose nodes are labelled by items such that, whenever a node is
labelled \(u\) and its children, in order, are labelled
\(u_1,\ldots,u_k\), the pair \((\langle u_1,\ldots,u_k\rangle,u)\) is in
\(\mathcal R\). An item is \emph{derivable} if it labels the root of some
derivation. Derivation from premises is the special case in which each
premise is added as an axiom. This is the format in which induction on
derivations follows the final rule, as \cref{sec:writing-foundations-04}
requires.

\begin{hbtheorem}{thm:p01-w02-p01-023-derivation-invariant}{proved=annex:p01-w02-p01-023-derivation-invariant}{Invariants of a rule system}
Let \(\mathcal R\) be a rule system on \(U\) and let \(P\subseteq U\)
satisfy: for every instance \((\langle u_1,\ldots,u_k\rangle,u)\in\mathcal R\),
if \(u_1,\ldots,u_k\in P\) then \(u\in P\). Then every derivable item lies in
\(P\). In particular, an item outside \(P\) is not derivable.
\end{hbtheorem}

A soundness theorem is an instance with \(P\) the set of valid items: the
proof of \cref{thm:p08-ipl-completeness-soundness} checks that every rule of
\(\mathsf{LJ}\) carries sequents valid in every Kripke model to a sequent
of the same kind. The second sentence of the theorem is the usual way of
proving that something is \emph{not} derivable: find a property that all
axioms have, that all rules preserve, and that the target item lacks.
\Cref{sec:p01-w02-p01-023-governing-logic} uses it in exactly this form.

A construction that runs through infinitely many stages and then takes their
union needs one more check, because a run has finite length.

\begin{hbtheorem}{thm:p01-w02-p01-023-limit}{proved=annex:p01-w02-p01-023-limit}{Invariants at a limit stage}
Let \(\langle S,I,\leadsto\rangle\) be a step system whose states are sets,
with \(s\subseteq t\) whenever \(s\leadsto t\). Let \(P\) be inductive and
let \(s_0,s_1,s_2,\ldots\) be an infinite sequence with \(s_0\in I\) and
\(s_n\leadsto s_{n+1}\) for every \(n\in\mathbb N\).
\begin{enumerate}
  \item Every \(s_n\) lies in \(P\).
  \item If \(P\) is \emph{closed under increasing unions}, that is,
        \(\bigcup_{n\in\mathbb N}t_n\in P\) whenever
        \(t_0\subseteq t_1\subseteq\cdots\) are members of \(P\), and if
        \(\bigcup_{n\in\mathbb N}s_n\in S\), then
        \(\bigcup_{n\in\mathbb N}s_n\in P\).
  \item The hypothesis of clause~2 cannot be dropped: in the step system
        whose states are the finite subsets of \(\mathbb N\) together with
        \(\mathbb N\), with \(I=\{\varnothing\}\) and
        \(s\leadsto s\cup\{\min(\mathbb N\setminus s)\}\) for finite \(s\),
        the property ``\(s\) is finite'' is inductive, the stages are
        \(s_n=\{0,\ldots,n-1\}\), and their union \(\mathbb N\) is not
        finite.
\end{enumerate}
\end{hbtheorem}

The Lindenbaum construction of \cref{lem:prop-lindenbaum} has this shape: the
invariant ``\(\Gamma_n\) is consistent'' is preserved at each stage, and the
union is consistent because consistency is closed under increasing unions,
a derivation of \(\bottom\) using only finitely many assumptions.

\begin{hbmethod}{def:p01-w02-p01-023-invariant-method}{Writing an invariant argument}
\begin{enumerate}
  \item \emph{Fix the step system.} State the states, the initial states and
        every kind of step with its enabling condition; for a calculus,
        state the items and every rule schema with its side conditions.
  \item \emph{State the invariant} \(P\) as a property of one state (or one
        item), with every parameter fixed. A property of the history of a
        run is not a property of a state; if the argument needs a fact
        about earlier states, add that fact, as it bears on the current
        state, to \(P\).
  \item \emph{Establishment.} Verify \(I\subseteq P\) (for a calculus: every
        axiom lies in \(P\)).
  \item \emph{Preservation.} For each kind of step separately, assume
        \(s\in P\) and the enabling condition of the step, and verify that
        the resulting state lies in \(P\), using no other information about
        \(s\). If a case needs more, strengthen \(P\) and repeat steps 3
        and 4 for the stronger property.
  \item \emph{Close.} Show \(P\subseteq Q\) for a statement about every
        reachable state, \(P\cap T\subseteq Q\) for a statement about the
        terminal states, or that the target item lies outside \(P\) for a
        non-derivability statement.
  \item \emph{Terminate, if the conclusion is about a final state that must
        exist.} Give a variant for \(P\) (\cref{thm:p01-w02-p01-023-variant}).
  \item \emph{Pass to a limit, if the construction takes a union of
        infinitely many stages.} Prove closure under increasing unions, or
        give another argument for the limit
        (\cref{thm:p01-w02-p01-023-limit}).
\end{enumerate}
\end{hbmethod}

\paragraph{Worked example: forward chaining for Horn clauses.}
\Cref{thm:p02-w02-p02-019-horn-procedure} proves correctness and the
linear bound for the clause-counter implementation, under the input and
cost conventions of \cref{def:p02-w02-p02-019-clause-input}.
The present example isolates the invariant and termination argument,
without identifying the number of abstract inference steps with the
implementation's running time. A Horn clause has at most one positive
literal (\cref{def:p02-tractable-classes}). Write a Horn
clause \(\neg b_1\lor\cdots\lor\neg b_k\lor h\) as the pair \((B,h)\) with
\(B=\{b_1,\ldots,b_k\}\) and \(h\) its positive literal, and write the
clause \(\neg b_1\lor\cdots\lor\neg b_k\) with no positive literal as
\((B,\bottom)\); the empty clause is \((\varnothing,\bottom)\). A valuation
\(v\) satisfies \((B,h)\) exactly when \(v\) makes some \(b\in B\) false,
or \(h\) is an atom that \(v\) makes true.

\begin{hbdefinition}{def:p01-w02-p01-023-forward-chaining}{Forward chaining as a step system}
Let \(\Phi\) be a finite set of Horn clauses in the pair form above, and let
\(A\) be the finite set of atoms occurring in \(\Phi\). The forward-chaining
step system of \(\Phi\) has as states the subsets \(M\subseteq
A\cup\{\bottom\}\), the single initial state \(\varnothing\), and the steps
\[
  M\leadsto M\cup\{h\}
  \qquad\text{whenever }(B,h)\in\Phi,\ B\subseteq M\text{ and }h\notin M.
\]
A state \(M\) is \emph{closed} if \(h\in M\) for every \((B,h)\in\Phi\)
with \(B\subseteq M\). The property \(P_\Phi\) consists of the states
\(M\) such that every valuation that satisfies every clause of \(\Phi\)
makes every member of \(M\) true, where \(\bottom\) is true under no
valuation. The function \(\mu_\Phi\) is
\(\mu_\Phi(M)=\lvert A\cup\{\bottom\}\rvert-\lvert M\rvert\).
\end{hbdefinition}

\begin{hbtheorem}{thm:p01-w02-p01-023-horn-chaining}{proved=annex:p01-w02-p01-023-horn-chaining}{Correctness of forward chaining by invariants}
In the forward-chaining step system of a finite set \(\Phi\) of Horn clauses:
\begin{enumerate}
  \item \(P_\Phi\) is inductive;
  \item \(\mu_\Phi\) is a variant for \(P_\Phi\), so every run has length at
        most \(\lvert A\rvert+1\) and some run ends at a terminal state;
  \item a state is terminal if and only if it is closed; and
  \item for every reachable terminal state \(M\): \(\bottom\in M\) if and
        only if no valuation satisfies every clause of \(\Phi\). When
        \(\bottom\notin M\), the valuation that makes exactly the atoms
        in \(M\) true satisfies every clause of \(\Phi\).
\end{enumerate}
\end{hbtheorem}

The two halves of clause~4 close the argument differently. If \(\bottom\) is
reached, the invariant alone gives unsatisfiability: a satisfying valuation
would make \(\bottom\) true. If \(\bottom\) is not reached, the invariant
gives nothing, and the conclusion comes from the exit condition, closure,
which produces a satisfying valuation.

\begin{hbexample}{def:p01-w02-p01-023-horn-example}{Two runs of forward chaining}
Let \(p,q,r\) be distinct atoms and
\(\Phi=\{(\varnothing,p),(\{p\},q),(\{p,q\},r),(\{q,r\},\bottom)\}\), the
clauses \(p\), \(\neg p\lor q\), \(\neg p\lor\neg q\lor r\) and
\(\neg q\lor\neg r\). One run is
\[
  \varnothing\leadsto\{p\}\leadsto\{p,q\}\leadsto\{p,q,r\}
  \leadsto\{p,q,r,\bottom\}.
\]
Each step uses the clause whose set \(B\) is contained in the current state
and whose \(h\) is new: \((\varnothing,p)\), then \((\{p\},q)\), then
\((\{p,q\},r)\), then \((\{q,r\},\bottom)\). The last state contains every
\(h\) of \(\Phi\), so it is closed and hence terminal. By
\cref{thm:p01-w02-p01-023-horn-chaining}, clause~4, \(\Phi\) is
unsatisfiable. For \(\Phi'=\Phi\setminus\{(\{q,r\},\bottom)\}\) the same
first three steps lead to the terminal state \(\{p,q,r\}\), which does not
contain \(\bottom\), and the valuation making \(p,q,r\) true satisfies
\(\Phi'\). The run length \(4\) for \(\Phi\) is within the bound
\(\lvert A\rvert+1=4\) of clause~2.
\end{hbexample}

\paragraph{Failure case: a true property that is not inductive.}
An invariant need not be inductive, and an argument that checks (I2) for a
property that is merely true of the reachable states fails, because (I2)
quantifies over every state with the property, reachable or not.

\begin{hbtheorem}{thm:p01-w02-p01-023-non-inductive}{proved=annex:p01-w02-p01-023-non-inductive}{An invariant that must be strengthened}
In the step system \(\langle\mathbb Z,\{0\},\leadsto\rangle\) with
\(s\leadsto t\) if and only if \(t=s+2\):
\begin{enumerate}
  \item \(\operatorname{Reach}=\{2m:m\in\mathbb N\}\);
  \item the property \(Q=\{s\in\mathbb Z:s\neq1\}\) is an invariant;
  \item \(Q\) is not inductive; and
  \item the property \(E\) of being even is inductive, and
        \(E\subseteq Q\).
\end{enumerate}
\end{hbtheorem}

\begin{hbwarning}{def:p01-w02-p01-023-invariant-limits}{What an invariant argument does not give}
An inductive property gives information about reachable states only. It
says nothing about termination, which needs a separate variant
(\cref{thm:p01-w02-p01-023-variant}; compare the separate total-correctness
while rule of \cref{sec:p09-program-logics}); nothing about a union of infinitely
many stages without a limit argument (\cref{thm:p01-w02-p01-023-limit},
clause~3); and nothing about states that are not reachable, so a closing
implication must be proved for every state of \(P\) (or of \(P\cap T\)),
not only for the states one has in mind. The variant of
\cref{thm:p01-w02-p01-023-variant} takes values in \(\mathbb N\); this
section proves nothing about variants with values in other orders.
\end{hbwarning}

\paragraph{Relationships.}
Structural induction, induction on derivations and induction on a complexity
measure (\cref{sec:writing-foundations-04}) are the special cases of
\cref{thm:p01-w02-p01-023-derivation-invariant} in which the rule system is
the grammar, the calculus or the measure ordering. A construction
(\cref{sec:writing-foundations-04}) is often verified by an invariant of its
stages, with \cref{thm:p01-w02-p01-023-limit} at the limit. Loop invariants
in the Hoare calculus are the iteration case
(\cref{def:p09-hoare-triple}, and the partial-correctness while rule of
\cref{sec:p09-program-logics}).

\subsection{Detecting a change of governing logic}
\label{sec:p01-w02-p01-023-governing-logic}

\paragraph{Purpose, inputs and conclusions.}
An argument has a governing logic: the logic in which its conclusion is
claimed to follow. A change of governing logic occurs when the argument,
while claiming a conclusion in one logic, uses a step that is valid only in
a stronger one, or applies a metatheorem proved for a different logic. The
remarks on contraposition and proof by contradiction in
\cref{sec:writing-foundations-04} identify two such steps for
intuitionistic logic. This subsection turns them into a check. Its inputs are
an argument, the stated object logic and the stated metalogic. Its
conclusions are, for each step, whether it is licensed in the logic of its
level, a semantic confirmation when it is not, and a repair.

The comparison is between intuitionistic propositional logic
\(\mathsf{IPL}\) and classical propositional consequence on the same
language. Formulas are built from the atoms in a set \(\mathsf{Prop}\) and
\(\bottom\) by \(\land,\lor,\to\), and \(\neg\varphi\) abbreviates
\(\varphi\to\bottom\), as in \cref{sec:intuitionistic-logic-01}.
\(\Gamma\derives_{\mathsf{IPL}}\varphi\) is finite-support derivability in
\(\mathsf{LJ}\), and \(\Gamma\entails_{\mathsf{IPL}}\varphi\) is Kripke
consequence at a common world, both as defined there and in
\cref{def:p08-ipl-completeness-semantics}. Write
\(\Gamma\entails_{\mathsf{CPL}}\varphi\) when every two-valued valuation
satisfying every member of \(\Gamma\) satisfies \(\varphi\), with the
truth clauses of \cref{sec:prop-syntax-semantics-02}; under every valuation
\(\varphi\to\bottom\) has the truth value of \(\neg\varphi\), so reading
negation as an abbreviation changes no classical truth value.

\begin{hbdefinition}{def:p01-w02-p01-023-governing-logic}{Governing logic, licensed step, change of governing logic}
Let an argument be given whose conclusion is a claim
\(\Gamma\derives_{L}\varphi\) in a stated object logic \(L\), and whose
reasoning about derivations, models or other objects is carried out in a
stated metalogic.
\begin{enumerate}
  \item An \emph{object-level step} is an inference inside the claimed
        \(L\)-argument: from judgements
        \(\Gamma_1\derives\varphi_1,\ldots,\Gamma_k\derives\varphi_k\) (the
        formulas available, with their open assumptions) to a judgement
        \(\Gamma_0\derives\varphi_0\). A \emph{meta-level step} is an
        inference about derivations, structures or other objects.
  \item An object-level step is \emph{licensed in \(L\)} if
        \(\Gamma_i\derives_{L}\varphi_i\) for every \(i=1,\ldots,k\) implies
        \(\Gamma_0\derives_{L}\varphi_0\). A meta-level step is licensed if
        it is valid in the stated metalogic.
  \item The argument \emph{changes its governing logic} if some
        object-level step is not licensed in \(L\), or if a meta-level step
        applies a theorem that was proved for a logic other than \(L\) to
        \(L\).
\end{enumerate}
\end{hbdefinition}

A step with no premise judgements, such as writing down \(p\lor\neg p\) with
no assumptions, is licensed in \(L\) exactly when its formula is
\(L\)-derivable from the available assumptions. The following theorem lists
classical principles whose use in an \(\mathsf{IPL}\)-argument is a change of
governing logic. Clause~1(c) is the instance of the contraposition schema
\((\neg\psi\to\neg\varphi)\to(\varphi\to\psi)\) with
\(\varphi:=\neg\neg p\) and \(\psi:=p\); it makes precise the remark in
\cref{sec:writing-foundations-04} that this schema is not generally
derivable intuitionistically.

\begin{hbtheorem}{thm:p01-w02-p01-023-classical-principles}{proved=annex:p01-w02-p01-023-classical-principles}{Classical principles that are not intuitionistic}
Let \(p\) and \(q\) be distinct atoms.
\begin{enumerate}
  \item Each of the following formulas \(\chi\)
        \begin{enumerate}
          \item[(a)] \(p\lor\neg p\) \hfill(excluded middle)
          \item[(b)] \(\neg\neg p\to p\) \hfill(double-negation elimination)
          \item[(c)] \((\neg p\to\neg\neg\neg p)\to(\neg\neg p\to p)\)
                \hfill(an instance of contraposition)
          \item[(d)] \(((p\to q)\to p)\to p\) \hfill(Peirce's formula)
          \item[(e)] \(\neg(p\land q)\to(\neg p\lor\neg q)\)
                \hfill(a De Morgan direction)
        \end{enumerate}
        satisfies \(\entails_{\mathsf{CPL}}\chi\) and
        \(\notderives_{\mathsf{IPL}}\chi\).
  \item \(\{\neg\neg p,\neg p\}\derives_{\mathsf{IPL}}\bottom\) and
        \(\{\neg\neg p\}\entails_{\mathsf{CPL}}p\), but
        \(\{\neg\neg p\}\notderives_{\mathsf{IPL}}p\). Hence the rule
        ``from \(\Gamma\cup\{\neg\varphi\}\derives\bottom\) infer
        \(\Gamma\derives\varphi\)'' is not licensed in \(\mathsf{IPL}\).
  \item \(\entails_{\mathsf{CPL}}p\lor\neg p\), but neither
        \(\entails_{\mathsf{CPL}}p\) nor \(\entails_{\mathsf{CPL}}\neg p\).
\end{enumerate}
\end{hbtheorem}

Clause~3 concerns the opposite direction. The disjunction property
(\cref{thm:disjunction-property}) is a theorem about \(\mathsf{LJ}\); an
argument that applies it to classical validity changes its governing logic
at the meta level, and clause~3 shows that the conclusion it would give is
false.

Not every step that resembles a classical principle is one. The next theorem
lists the intuitionistically sound forms, and the local hypotheses under
which the classical forms become sound.

\begin{hbtheorem}{thm:p01-w02-p01-023-repairs}{proved=annex:p01-w02-p01-023-repairs}{Intuitionistic replacements and local repairs}
For all formulas \(\varphi,\psi\), every set \(\Gamma\) of formulas, and
distinct atoms \(p,q\):
\begin{enumerate}
  \item[(a)] if \(\Gamma\cup\{\varphi\}\derives_{\mathsf{IPL}}\bottom\), then
        \(\Gamma\derives_{\mathsf{IPL}}\neg\varphi\) (proof of a negation);
  \item[(b)] if \(\Gamma\cup\{\neg\varphi\}\derives_{\mathsf{IPL}}\bottom\),
        then \(\Gamma\derives_{\mathsf{IPL}}\neg\neg\varphi\) (reductio with a
        weakened conclusion);
  \item[(c)] \(\derives_{\mathsf{IPL}}(\varphi\to\psi)\to(\neg\psi\to\neg\varphi)\)
        (contraposition in the forward direction);
  \item[(d)] \(\{\varphi\lor\neg\varphi\}\derives_{\mathsf{IPL}}\neg\neg\varphi\to\varphi\)
        and
        \(\{\psi\lor\neg\psi\}\derives_{\mathsf{IPL}}(\neg\psi\to\neg\varphi)\to(\varphi\to\psi)\);
  \item[(e)] \(\{p\lor\neg p\}\derives_{\mathsf{IPL}}\neg(p\land q)\to(\neg p\lor\neg q)\);
  \item[(f)] \(\derives_{\mathsf{IPL}}\neg\neg\bigl(\neg(p\land q)\to(\neg p\lor\neg q)\bigr)\).
\end{enumerate}
\end{hbtheorem}

\begin{hbmethod}{def:p01-w02-p01-023-logic-diagnostic}{Checking the governing logic of an argument}
\begin{enumerate}
  \item \emph{Record both logics.} Write down the object logic \(L\) of the
        conclusion and the metalogic of the reasoning. If either is unstated
        and the argument uses a principle of
        \cref{thm:p01-w02-p01-023-classical-principles}, the argument is
        incomplete until it is stated (\cref{sec:levels-01}).
  \item \emph{Assign each step a level.} A step inside the claimed
        derivation is object-level; a step about derivations, models or
        valuations is meta-level.
  \item \emph{Check each object-level step against \(L\).} Name the rule or
        derived rule that licenses it. When \(L\) is \(\mathsf{IPL}\), flag
        every step whose only justification is one of the forms of
        \cref{thm:p01-w02-p01-023-classical-principles}: a case split on
        \(\varphi\lor\neg\varphi\) when that disjunction has been neither
        derived nor assumed,
        passing from \(\neg\neg\varphi\) to \(\varphi\), concluding
        \(\varphi\) (rather than \(\neg\varphi\)) from a contradiction,
        passing from \(\neg\psi\to\neg\varphi\) to \(\varphi\to\psi\),
        Peirce's formula, or passing from \(\neg(\varphi\land\psi)\) to
        \(\neg\varphi\lor\neg\psi\). Do not flag the forms of
        \cref{thm:p01-w02-p01-023-repairs}(a)--(c).
  \item \emph{Check each meta-level step against the metalogic.} Classical
        case analysis about Kripke models or derivations is licensed in a
        classical metalogic even when the object logic is intuitionistic;
        the proofs of
        \cref{thm:p01-w02-p01-023-classical-principles,thm:p01-w02-p01-023-repairs}
        are of this kind.
  \item \emph{Check each appeal to a metatheorem} (soundness, completeness,
        the disjunction property, Glivenko's theorem) against the logic for
        which it was proved.
  \item \emph{Confirm a flag semantically.} A flagged step may still be
        licensed by some other derivation. It is not licensed in
        \(\mathsf{IPL}\) if a Kripke model has a world forcing its premises
        but not its conclusion (\cref{thm:p08-ipl-completeness-soundness}).
  \item \emph{Repair}, and say which repair is used:
        \begin{enumerate}
          \item[R1] reprove the step in \(L\) without the principle, for
                example by \cref{thm:p01-w02-p01-023-repairs}(a) when the
                target is a negation;
          \item[R2] weaken the conclusion, for example to \(\neg\neg\varphi\)
                by \cref{thm:p01-w02-p01-023-repairs}(b), or by Glivenko's
                theorem (\cref{thm:glivenko}) for a propositional classical
                derivation;
          \item[R3] add the needed instance, such as
                \(\varphi\lor\neg\varphi\), as an explicit hypothesis and
                carry it into the statement of the result
                (\cref{thm:p01-w02-p01-023-repairs}(d), (e));
          \item[R4] restate the result for the stronger logic, and relabel it
                (the repair of a strength error in
                \cref{sec:writing-foundations-07}).
        \end{enumerate}
\end{enumerate}
\end{hbmethod}

\begin{table}[tp]
\centering
\small
\begin{tabular}{@{}p{0.2\textwidth}p{0.36\textwidth}p{0.36\textwidth}@{}}
\toprule
Error & Diagnostic question & Repair \\
\midrule
Governing-logic error & Is every object-level step licensed in the stated
object logic, every meta-level step in the stated metalogic, and every cited
metatheorem proved for the logic it is applied to? & Name the rule used;
reprove without the principle, weaken the conclusion, add the instance as a
hypothesis, or restate the result for the stronger logic. \\
\bottomrule
\end{tabular}
\caption{The governing-logic entry, in the format of the error taxonomy of
\cref{sec:writing-foundations-07}.}
\label{tab:p01-w02-p01-023-logic-error}
\end{table}

\begin{hbexample}{def:p01-w02-p01-023-de-morgan-example}{A hidden case split on excluded middle}
Claim: \(\derives_{\mathsf{IPL}}\neg(p\land q)\to(\neg p\lor\neg q)\) for
distinct atoms \(p,q\). Proposed argument:
\begin{enumerate}
  \item assume \(\neg(p\land q)\);
  \item either \(p\) or \(\neg p\);
  \item if \(\neg p\), then \(\neg p\lor\neg q\);
  \item if \(p\): assume \(q\); then \(p\land q\), which with step~1 gives
        \(\bottom\); so \(\neg q\), and \(\neg p\lor\neg q\);
  \item by the cases of step~2, \(\neg p\lor\neg q\); discharging step~1
        gives the claim.
\end{enumerate}
\emph{Diagnostic.} The object logic is \(\mathsf{IPL}\); every step is
object-level. Steps 1, 3 and 5 are assumption, disjunction introduction,
disjunction elimination and implication introduction, none of which is a
flagged form. Step~4 derives
\(\neg q\) from a contradiction under the assumption \(q\): this is the proof
of a negation, licensed by \cref{thm:p01-w02-p01-023-repairs}(a). Step~2
has no premise judgement and asserts \(p\lor\neg p\), which is not
\(\mathsf{IPL}\)-derivable
(\cref{thm:p01-w02-p01-023-classical-principles}, clause~1(a)); it is
flagged.

\emph{Confirmation.} The flagged step is essential: by clause~1(e) of the
same theorem, the claim itself is not \(\mathsf{IPL}\)-derivable. The
countermodel in its proof has a root \(w_0\) with two incomparable
successors, \(p\) forced only at one and \(q\) only at the other; the root
forces \(\neg(p\land q)\) but neither \(\neg p\) nor \(\neg q\).

\emph{Repairs.} R3: with step~2 as a hypothesis the argument is an
\(\mathsf{IPL}\)-argument, and its conclusion is
\(\{p\lor\neg p\}\derives_{\mathsf{IPL}}\neg(p\land q)\to(\neg p\lor\neg q)\)
(\cref{thm:p01-w02-p01-023-repairs}(e)). R2: without hypotheses,
\(\derives_{\mathsf{IPL}}\neg\neg(\neg(p\land q)\to(\neg p\lor\neg q))\)
(\cref{thm:p01-w02-p01-023-repairs}(f)). R4: the original claim holds for
classical consequence (\cref{thm:p01-w02-p01-023-classical-principles},
clause~1(e)).
\end{hbexample}

\paragraph{Limitations.}
The list of flagged forms is the list proved in
\cref{thm:p01-w02-p01-023-classical-principles}; it is not a complete list of
classical principles, and a step outside it is checked by finding a
derivation or a countermodel. The diagnostic is stated for propositional
\(\mathsf{IPL}\) against classical consequence. For other pairs of logics
the definition of a licensed step applies unchanged, but the flagged forms
and the confirming semantics must be those of the logics concerned; this
section proves nothing about them. Clause~1 of
\cref{thm:p01-w02-p01-023-classical-principles} shows that particular
instances fail; it does not decide which other instances of the same schemas
are derivable.

\paragraph{Relationships.}
A change of governing logic is related to the level error of
\cref{sec:writing-foundations-07}, since the remedy begins by stating the
logic at each level, and to the strength error, since repair R4 is its
repair. The semantic confirmation is an application of
\cref{thm:p01-w02-p01-023-derivation-invariant}: Kripke validity is preserved
by every rule of \(\mathsf{LJ}\), and the flagged formulas lack it. The
negative translations of \cref{sec:intuitionistic-classical} systematize
repair R2.

\paragraph{Sources.}
The remarks on the ambient logic of contraposition and contradiction and the
error taxonomy are inherited from the frozen manuscript
\cite[lines 1505--1553, 1592--1631]{FrozenFormalLogicProofMethods}; the
Kripke semantics, persistence and the two-world counterexample used in the
proofs are its intuitionistic chapter
\cite[lines 12989--13127]{FrozenFormalLogicProofMethods}. The program
verification formulations are cited above to their pages. The definitions,
theorems and proofs of this section are editorial and are not attributed to
these sources.
